% Part: proof-theory
% Chapter: sequent-calculus
% Section: introduction

\documentclass[../../../include/open-logic-section]{subfiles}

\begin{document}

\olfileid{pt}{seq}{int}

\olsection{مقدمه}

حساب‌های دنباله‌ای دستگاه‌های !!{proof} هستند که گرهارد گنتسن در دههٔ
۱۹۳۰ معرفی کرد. هدف اصلی از آن‌ها فراهم‌کردن روشی عملی برای !!{proving}
نبود. گنتسن آن‌ها را به کار گرفت تا بتواند نتایجی دربارهٔ ساختارهای ممکنِ
!!{proof}s و ویژگی‌های آن‌ها اثبات کند. این‌گونه پرسش‌ها دقیقاً موضوع
نظریهٔ برهان‌اند.

حساب‌های دنباله‌ای، چنان‌که نامشان نشان می‌دهد، بر \emph{دنباله‌ها}
عمل می‌کنند، نه بر~!!{formula}s. هر دنباله جفتی از رشته‌ها یا چندمجموعه‌های
!!{formula}s است. دستگاه‌های اصلی گنتسن (~\Log{LK} و~\Log{LJ}) از
رشته‌ها استفاده می‌کردند؛ نظریه‌پردازان برهان امروزه چندمجموعه‌ها را ترجیح
می‌دهند. چندمجموعه مانند مجموعه است، با این تفاوت که می‌تواند !!{element}
را بیش از یک بار در بر داشته باشد. بااین‌حال، همانند مجموعه‌ها، ترتیب
!!{element}s اهمیتی ندارد. پس عبارت‌های~$\{!A, !A, !B\}$ و
$\{!A, !B, !A\}$ یک چندمجموعه را وصف می‌کنند، اما آن چندمجموعه با
$\{!A, !B\}$ و~$\{!A, !A, !B, !B\}$ متفاوت است، زیرا تعداد رخدادهای
$!A$ و~$!B$ در آن‌ها متفاوت است.

در نظریهٔ برهان، چندمجموعه‌های !!{formula}s را به‌طور سنتی با حروف بزرگ
یونانی می‌نویسند. پس~$\Gamma$، $\Delta$، $\Pi$، $\Lambda$ یا~$\Theta$
را چندمجموعه‌هایی از !!{formula}s منطق مورد بحث می‌گیریم. همچنین به‌جای
نوشتن~$\tuple{\Gamma, \Delta}$، دو مؤلفه را به‌طور سنتی با نماد دنباله
جدا می‌کنند؛ ما از~$\Sequent$ استفاده می‌کنیم.

\begin{defn}[دنباله]
یک \emph{دنباله} عبارتی به صورت
\[
\Gamma \Sequent \Delta
\]
است که در آن~$\Gamma$ و~$\Delta$ چندمجموعه‌های متناهی (و احتمالاً تهی)
از~!!{formula}s هستند. $\Gamma$ را \emph{مقدم} و~$\Delta$ را
\emph{تالی} می‌نامند.
\end{defn}

فرض کنید~$!G$ عطفِ !!{formula}s موجود در~$\Gamma$ باشد (و عطف تهی را
$\ltrue$ بگیریم) و~$!D$ فصلِ !!{formula}s موجود در~$\Delta$ باشد (و فصل
تهی را~$\lfalse$ بگیریم). آنگاه دنبالهٔ~$\Gamma \Sequent \Delta$ در
ساختاری چون~$\Struct{M}$ صادق است اگر و تنها اگر~$!G \lif !D$
صادق باشد. در منطق کلاسیک، این همان است که یا یکی از !!{formula}s موجود
در~$\Gamma$ کاذب باشد یا یکی از~!!{formula}s موجود در~$\Delta$ صادق.

اگر~$\Gamma$ چندمجموعه‌ای از !!{formula}s باشد، $\Gamma, !A$ (یا
$!A, \Gamma$) را برای حاصل افزودن~$!A$ به~$\Gamma$ می‌نویسیم. اگر
$\Delta$ نیز چندمجموعه‌ای از !!{formula}s باشد، $\Gamma, \Delta$ اجتماع
چندمجموعه‌ایِ آن دو است: اگر~$\Gamma$، $!A$ را با
\emph{چندگانگی}~$n$ در بر داشته باشد (یعنی~$n$ نسخه از~$!A$ داشته
باشد) و~$\Delta$، $!A$ را با چندگانگی~$m$ در بر داشته باشد، آنگاه
$\Gamma, \Delta$، $!A$ را با چندگانگی~$n+m$ در بر دارد. اگر چندمجموعهٔ
یک سوی دنباله تهی باشد، معمولاً جای آن را خالی می‌گذاریم. برای نمونه،
$\Sequent !A$ دنباله‌ای با مقدم تهی و تالی‌ای است که~$!A$ را با چندگانگی
$1$ در بر دارد.

!!^a{proof} در حساب دنباله‌ای درختی از دنباله‌هاست که دنباله‌های بالایی
(یا آغازی) آن اصول دستگاه مورد نظرند، و هر دنباله‌ای که زیر یک یا دو
دنبالهٔ دیگر قرار گیرد باید به‌درستی با قاعده‌ای استنتاجی در دستگاه از آن‌ها
به دست آید. نخست دو دستگاه دنباله‌ای متفاوت برای منطق کلاسیک،
$\Log{G1c}$ و~$\Log{G3c}$، را بررسی می‌کنیم. اصول و قواعد آن‌ها در
\cref{pt:seq:int:tab:G1c,pt:seq:int:tab:G3c} آمده‌اند.
\oliflabeldef{fol:seq::chap}{ (دستگاه اصلی گنتسن~\Log{LK} در
\olref[fol][seq][]{chap} شرح داده شده است.)}{}

\subfile{rules-G1c}
\subfile{rules-G3c}

دو دستگاه تفاوت‌هایی ظریف دارند و همین تفاوت‌ها ویژگی‌های نظریه‌برهانی
آن‌ها را متفاوت می‌کند. اصول~\Log{G1c} به صورت~$!A \Sequent !A$ هستند
و هیچ !!{formula} دیگری در آن‌ها مجاز نیست. اصول~\Log{G3c} به صورت
$!C, \Gamma \Sequent \Delta, !C$ هستند و~$\Gamma$ و~$\Delta$ می‌توانند
!!{formula}s «جانبی» دلخواه داشته باشند، اما !!{formula}~$!C$ که در هر
دو سو می‌آید باید \emph{اتمی} باشد. \Log{G1c} از قواعد
$\LeftR{\land}$ و~$\RightR{\lor}$ هر یک دو نسخه دارد، حال آنکه
\Log{G3c} از هر یک تنها یک نسخه دارد. همچنین~\Log{G1c} قواعد ساختاری
اضافیِ «انقباض» (~\LeftR{\Contraction}، \RightR{\Contraction}) و
«تضعیف» (~\LeftR{\Weakening}، \RightR{\Weakening}) را دارد.

اتفاقاً هر دو دستگاه~\Log{G1c} و~\Log{G3c} برای منطق کلاسیک درست و کامل
هستند. به بیان دیگر، هر دنبالهٔ !!{provable} در این دستگاه‌ها در هر
!!{structure} صادق است، و هر دنباله‌ای که در هر !!{structure} صادق باشد
!!a{proof} دارد. پس پرسش از اینکه \emph{آیا} دنباله‌ای !!a{proof}
دارد، با روش‌های دیگر منطق، برای مثال روش‌های نظریهٔ مدل، نیز پاسخ‌پذیر است.
و چون هر دو دستگاه دقیقاً دنباله‌هایی را !!{prove} می‌کنند که در هر
!!{structure} صادق‌اند، به‌ویژه همان دنباله‌ها را اثبات می‌کنند.

بااین‌حال، نظریهٔ برهان تنها به این نمی‌پردازد که \emph{آیا} چیزی
اثبات‌پذیر است، بلکه می‌پرسد \emph{چگونه}. نظریه‌پردازان برهان
پرسش‌هایی از این دست را بررسی می‌کنند: «آیا همیشه می‌توان از قاعده‌ای معین
پرهیز کرد؟»، «آیا می‌توان این قاعده را بی‌آنکه دنباله‌های تازه‌ای اثبات‌پذیر
شوند به دستگاه افزود؟» یا «آیا می‌توان !!{proof}s یک دستگاه را به
!!{proof}s دستگاهی دیگر تبدیل کرد؟» اگر پاسخ چنین پرسشی مثبت باشد، اثبات
آن غالباً با استقرا بر پیچیدگی !!{proof}s، یا هم‌ارز آن، با استقرا بر
ساختار !!{proof}s انجام می‌شود. این کار نه‌فقط برهانی برای واقعیت مورد نظر
(مثلاً اینکه اگر~\Log{G1c} دنباله‌ای را اثبات کند،~\Log{G3c} نیز همان را
اثبات می‌کند)، بلکه روالی نیز به دست می‌دهد (مثلاً برای تبدیل !!{proof}s
در~\Log{G1c} به~!!{proof}s در~\Log{G3c}).

یکی از نتایج مرکزی نظریهٔ برهان \emph{قضیهٔ حذف برش} است. این قضیه
نشان می‌دهد می‌توان قاعدهٔ برش
\[
\Axiom$\Gamma \fCenter \Delta, !A$
\Axiom$!A, \Pi \fCenter \Lambda$
\RightLabel{\Cut}
\BinaryInf$\Gamma, \Pi \fCenter \Delta, \Lambda$
\DisplayProof
\]
را بی‌آنکه دنباله‌های !!{provable} تغییر کنند به دستگاه‌های دنباله‌ای افزود.
افزون بر این، برهان قضیه روشی به دست می‌دهد که هر !!{proof} با قواعد
\Log{G1c} (یا~\Log{G3c}) همراه با~\Cut{} را به برهانی تبدیل می‌کند
که تنها قواعد~\Log{G1c} (یا~\Log{G3c}) را به کار می‌برد. به بیان دیگر،
این روش قاعدهٔ~\Cut{} را از !!{proof}sی که آن را به کار برده‌اند «حذف»
می‌کند؛ ازاین‌رو نام قضیه چنین است.

\end{document}
